\documentclass[11pt]{article}

\usepackage[margin=1in]{geometry}
\usepackage[T1]{fontenc}
\usepackage{lmodern}
\usepackage{amsmath}
\usepackage{amssymb}
\usepackage{graphicx}
\usepackage{hyperref}

\hypersetup{colorlinks=true, linkcolor=blue, citecolor=blue, urlcolor=blue}

\newcommand{\CC}{\mathbb{C}}
\newcommand{\RR}{\mathbb{R}}
\newcommand{\HH}{\mathbb{H}}
\newcommand{\OO}{\mathbb{O}}
\newcommand{\rSM}{r_{\mathrm{SM}}}

\title{No-Go Results for a Three-Qubit Entropy Ansatz for Gauge-Coupling Hierarchies, with an Exact Octonionic Rotor Construction in an Appendix}

\author{William T. Trevena\\[3pt]
{\small\itshape Independent Researcher (PhD, ISE, University of Florida)}\\
{\small\texttt{trevenaw7@gmail.com}}}

\date{June 2026}

\begin{document}

\maketitle

{\small\itshape Acknowledgment: Portions of this manuscript were developed with the assistance of Claude (Anthropic), an AI language model, which contributed to literature synthesis, mathematical exposition, drafting, and the computational experiments. An independent verification pass --- which identified two computational errors and several overclaims corrected in this version, as well as the closed form in Appendix~\ref{sec:rotor} --- was performed by a separately instructed instance of the same AI system; the repository includes the independent verification log and a response-to-review document. Consistent with COPE authorship guidelines, the AI system is acknowledged as a tool rather than listed as a coauthor.}

\begin{abstract}
Motivated by Szangolies's three-qubit octonionic Hopf construction of the $SU(3) \times SU(2) \times U(1)/\mathbb{Z}_6$ gauge group, we test a natural but naive extension of it: that the measured three-gauge-coupling hierarchy could be encoded, under a two-parameter logarithmic entropy map, in the single-qubit entanglement entropies of a three-qubit candidate state. We are not refuting Szangolies's construction; we are refuting this coupling-entropy extension, with scoped no-go results verified by published code. First, underconstraint: the map collapses to one entropy-gap-ratio condition, $r_S = \rSM$, a single scalar condition lying strictly inside the image of the three-qubit marginal polytope, so the matching set is positive-dimensional and matching carries no evidence; a sampling check across $10^7$ Haar-random states confirms the matching manifold is statistically generic, and null controls (random coupling targets, four-qubit systems, relabeled qubits) confirm that parameter freedom drives matching. We exhibit an explicit analytic one-parameter family of W-class (zero 3-tangle) states satisfying the ratio condition exactly, so matching does not even fix the SLOCC class. Principal obstruction: recognizing Furey's $\mathrm{Cl}(6)$ as the three-qubit operator algebra (Jordan--Wigner) yields charge as Hamming-weight grading --- and shows that gauge transformations are not local unitaries of the qubit factorization. The naive entropy observable is not gauge-invariant, so the map is not physically meaningful as stated. Independently, a scale-independent entropy triple is inconsistent with renormalization-group running; the surviving affine entropy-flow map is only a consistency envelope (feasible for $|B| \gtrsim 126$, binding at $M_Z$ alone, GUT-normalization dependent). An appendix supplies a dynamical control: the algebra-canonical rotor Hamiltonian $i(L_u + 2R_u)$ selects a unique, convention-invariant, permutation-\emph{symmetric} ground state $(\sqrt{7}\,e_0 - iu)/\sqrt{14}$ (marginal spectra $(1/7, 6/7)$, 3-tangle $8/49$): even the most canonical dynamical octonionic construction produces symmetry, not hierarchy. We close with an explicit map of what is ruled out and what is not.
\end{abstract}

\section{Introduction}
\label{sec:intro}

\subsection{Motivation: Division Algebras, the Szangolies Bridge, and Entanglement}
\label{sec:motivation}

The four normed division algebras --- $\RR$, $\CC$, $\HH$, $\OO$ --- are the only algebras over the reals admitting a multiplicative norm (Hurwitz \cite{hurwitz}). Fifty years of work has pursued Standard Model structure in the largest of them: G\"unaydin and G\"ursey identified SU(3) color within the octonion automorphisms \cite{gunaydin}; Dixon proposed $\RR \otimes \CC \otimes \HH \otimes \OO$ as the algebraic design of one generation with family replication \cite{dixon94, dixon04}; Furey constructed $SU(3)_C \times SU(2)_L \times U(1)_Y$ from division-algebraic ladder operator symmetries in $\mathrm{Cl}(6)$, with correct electric charges and charge quantization from the integer spectrum of a number operator \cite{furey14, furey18}; Todorov and Dubois-Violette derived the same symmetry from the exceptional Jordan algebra \cite{todorov18, todorov22}.

Szangolies (2025) connected this program to quantum information \cite{szangolies}. Building on Mosseri--Dandoloff \cite{mosseri} and Bernevig--Chen \cite{bernevig}, he showed that entangled multi-qubit systems map to Hopf fibrations over the normed division algebras:

\begin{center}
\begin{tabular}{cccc}
\hline
Qubits & Algebra & Hopf Fibration & Residual Symmetry \\
\hline
1 & $\CC$ & $S^1 \to S^3 \to S^2$ & $U(1)$ \\
2 & $\HH$ & $S^3 \to S^7 \to S^4$ & $SU(2) \times U(1)$ \\
3 & $\OO$ & $S^7 \to S^{15} \to S^8$ & $SU(3) \times SU(2) \times U(1)/\mathbb{Z}_6$ \\
\hline
\end{tabular}
\end{center}

The three-qubit octonionic case is the unique normed-division-algebra Hopf fibration producing the Standard Model gauge group, and the Hopf maps are entanglement-sensitive: separable states produce no non-abelian gauge structure \cite{szangolies}. A second, independent program --- emergent spacetime from entanglement, from the Bekenstein area law through AdS/CFT and Ryu--Takayanagi to the Page--Wootters mechanism \cite{bekenstein, maldacena, ryu, page, takayanagi} --- supplies the broader motivation for treating entanglement as physically primary. We note that the nexus of three qubits, the Cayley hyperdeterminant, the octonions, and the Fano plane is also the subject of an established independent literature, the black-hole/qubit correspondence \cite{borsten}; our use of the nexus (coupling hierarchies and state selection) is, to our knowledge, distinct from its use there.

These motivations jointly suggest a specific quantitative question, the subject of this paper: \emph{is the measured three-gauge-coupling hierarchy encoded, under a proposed three-qubit entropy map, in the entanglement entropies of a three-qubit candidate state?} The body of the paper answers no, in every form of the question we could make precise, and the central finding is that the question itself is ill-posed until a gauge-invariant entanglement functional is specified. One point of scope deserves emphasis at the outset: \textbf{we are not refuting Szangolies's construction}, which concerns gauge-\emph{group} structure and makes no claim about coupling values; \textbf{we are refuting a natural but naive coupling-entropy extension of it} --- the idea that, since the construction distinguishes three tensor factors associated with the three gauge factors, the three measured couplings might be read off from the three single-qubit entropies.

\subsection{Summary of Results}
\label{sec:summary}

The paper is organized as a sequence of no-go results of increasing depth, closed by an exact dynamical control in an appendix.

\begin{enumerate}

\item \textbf{Underconstraint} (Section~\ref{sec:static}). Under the two-parameter logarithmic map, the coupling fit collapses to a single entropy-gap-ratio condition, $r_S = \rSM \approx 1.82$ (Lemma~1). The core point is geometric: this is one scalar condition imposed inside a feasible region of entropy triples, so the matching set is positive-dimensional (Proposition~1) --- a matching state always exists, and matching is not evidence. Two-qubit states cannot encode unequal couplings at all ($S_1 = S_2$ identically), so three qubits is the smallest system on which the ansatz has any content.

\item \textbf{A sampling sanity check} (Section~\ref{sec:generic}). A check across $10^7$ Haar-random states confirms the analytic picture: the matching manifold is statistically generic in its entanglement invariants (Kolmogorov--Smirnov $D = 0.031$ against ordering-conditioned controls), and four null controls --- random coupling targets, two- and four-qubit systems, relabeled qubits --- confirm that parameter freedom, not algebraic structure, drives matching.

\item \textbf{The W-class counterexample} (Section~\ref{sec:wclass}). We construct an explicit analytic one-parameter family of weighted-W states ($\tau_3 = 0$) satisfying $r_S = \rSM$ exactly (Lemma~2, Section~\ref{sec:wfamily}), upgrading the numerical matching curve to a closed-form parameterization: coupling matching does not even fix the SLOCC class. By contrast, every canonical octonionic state construction considered here is permutation-degenerate and cannot produce a hierarchy at all.

\item \textbf{The gauge-covariance obstruction} (Section~\ref{sec:jw}). Furey's $\mathrm{Cl}(6)$ construction is the three-qubit operator algebra under the Jordan--Wigner transformation; electric charge becomes Hamming-weight grading. The principal yield is an obstruction: gauge transformations are not local unitaries of the qubit factorization, so single-qubit entropies are gauge-frame dependent. The naive entropy observable is not gauge-invariant; therefore the map is not physically meaningful. This is the sharpest result of the paper, and the conclusion the rest of the paper is organized around.

\item \textbf{The renormalization-group obstruction} (Section~\ref{sec:rg}). Independently of the gauge obstruction, the gap ratio $\rSM(\mu)$ runs ($\sim$24\% between $M_Z$ and $10^6$ GeV, with a pole near $10^{13}$ GeV), so a scale-independent entropy triple cannot match the couplings at all scales. This is true but not deep; we keep the section short. Its one retained technical content is the entropy-flow \emph{consistency envelope}: feasibility forces $|B| \gtrsim 126$, the constraint binds at $M_Z$ alone, no observable is constrained, and the single falsifiable consequence (symmetric boundary state $\Leftrightarrow$ exact unification) restates the unification question and inherits the GUT normalization convention.

\item \textbf{A dynamical control} (Appendix~\ref{sec:rotor}). The rotor Hamiltonian with multiplication-table-sign couplings equals $i(L_u + 2R_u)$, $u = \sum_a e_a$; its unique ground state $(\sqrt{7}\,e_0 - iu)/\sqrt{14}$ is provably convention-invariant and permutation-symmetric, with exact invariants (marginal spectra $(1/7, 6/7)$, 3-tangle $= 8/49$). This is the one place in the paper where ``vacuum'' is used literally --- the algebra-canonical dynamics selects the state --- and it functions as a control on the whole program: even the most canonical dynamical octonionic construction produces symmetry, not hierarchy.

\end{enumerate}

Section~\ref{sec:discussion} closes with a table of what is ruled out and what is not, and reproducibility information.

\subsection{Scope, Terminology, and Limitations}
\label{sec:scope}

\textbf{Scope of the no-go results.} Every negative result in this paper is scoped: it applies to the specific ansatz stated (single-qubit von Neumann entropies of a three-qubit pure state, mapped to couplings by the two-parameter logarithmic map of Section~\ref{sec:conjecture}, with the qubit-gauge assignment fixed by the Szangolies construction), and to the specific families of states and Hamiltonians enumerated where each result is proved. We rule out the naive static matching program in this scope; we do not rule out gauge-invariant or algebra-relative formulations (see Table~\ref{tab:notruledout}).

\textbf{Terminology.} The object we test is not the Standard Model --- it has no gauge fields, matter, chirality, anomalies, Yukawas, Higgsing, thresholds, or scheme dependence. Where precision matters we therefore say ``the measured three-gauge-coupling hierarchy under a proposed three-qubit entropy map,'' not ``the Standard Model coupling hierarchy.'' Likewise, the three-qubit state under test is a \emph{candidate state} (a minimal-state ansatz in $\CC^2 \otimes \CC^2 \otimes \CC^2$), not a QFT vacuum; we reserve the word \emph{vacuum} for Appendix~\ref{sec:rotor}, where a Hamiltonian genuinely selects its ground state, and we say so explicitly there.

\textbf{Three distinct entropies.} Discussions in this area move between three inequivalent notions of ``entanglement entropy,'' and blurring them is a standing source of confusion. Table~\ref{tab:notions} fixes the distinction once; this paper's no-go results concern notion (i) only.

\begin{table}[ht]
\centering
\small
\begin{tabular}{p{3.4cm}p{5.6cm}p{4.6cm}}
\hline
\textbf{Notion} & \textbf{Definition} & \textbf{Role in this paper} \\
\hline
(i) Subsystem entropy & von Neumann entropy $S_i = -\mathrm{tr}\,\rho_i \log_2 \rho_i$ of a single qubit's reduced density matrix, relative to a fixed tensor factorization & \textbf{The object under test.} Shown gauge-frame dependent in Section~\ref{sec:jw} \\
(ii) Algebraic entropy & Entropy relative to a gauge-invariant operator subalgebra (commutant of the gauge action), in the sense of Barnum--Knill--Ortiz--Viola \cite{barnum} and Zanardi \cite{zanardi} & Not computed here; the gauge-invariant repair, carried out in follow-up work (Note Added) \\
(iii) Flux-sector (edge) entropy & Shannon entropy of the superselection-sector (flux/charge) distribution across an entangling cut in gauge theory & Not used here; the notion relevant to lattice-gauge edge modes \\
\hline
\end{tabular}
\caption{Three inequivalent notions of entanglement entropy. The no-go results of this paper concern notion (i), the naive subsystem entropy; its failure of gauge invariance (Section~\ref{sec:jw}) is precisely what motivates notion (ii).}
\label{tab:notions}
\end{table}

\textbf{Heuristic, not ontology.} Earlier versions of this program organized the motivations of Section~\ref{sec:motivation} under an information-first framework. None of the technical results below, positive or negative, depends on accepting any such reading, and we do not develop it here; a short paragraph in Section~\ref{sec:broader} records the broader topics we deliberately set aside. We do not claim to derive the Standard Model, resolve quantum gravity, prove the number of fermion generations, or solve the cosmological constant problem.

\section{The Static Conjecture and Its Underconstraint}
\label{sec:static}

\subsection{The Measured Couplings and the Proposed Map}
\label{sec:conjecture}

The three gauge couplings at $M_Z \approx 91.2$ GeV \cite{pdg}:
\begin{align*}
\alpha_3 &= 0.1179 \pm 0.0010 && (SU(3)) \\
\alpha_2 &= 0.03374 \pm 0.00003 && (SU(2)) \\
\alpha_1 &= 0.01695 \pm 0.00002 && (U(1),\ \text{GUT-normalized: } \alpha_1 = \tfrac{5}{3}\alpha_Y)
\end{align*}

These are free parameters --- measured, not derived from any known principle.

\textbf{Normalization caveat (stated at the outset, because every $\alpha_1$-dependent number below inherits it).} The abelian coupling enters this paper in the GUT normalization $\alpha_1 = (5/3)\alpha_Y$. That factor is a convention inherited from embedding hypercharge in SU(5)-type groups; nothing in the three-qubit framework fixes it. Under a different admissible normalization, the target ratio $\rSM$, the coupling-crossing scales, and the unification asymmetries of Section~\ref{sec:envelope} all change. Any framework that declines GUT embedding owes an independent derivation of the abelian normalization; none currently exists (Section~\ref{sec:corollary}).

\textbf{Conjecture (Logarithmic Coupling-Entropy Map).} There exists a three-qubit pure candidate state $|\Omega\rangle \in \CC^2 \otimes \CC^2 \otimes \CC^2$ and real constants $A$, $B$ such that:
\begin{equation}
\log(\alpha_i^{-1}) = A + B \cdot S_i \qquad \text{for } i = 1, 2, 3
\label{eq:logmap}
\end{equation}
where $S_i$ is the von Neumann entropy of the $i$-th qubit's reduced density matrix, and the qubit-gauge assignment (qubit 1 $\leftrightarrow$ $SU(3)$, qubit 2 $\leftrightarrow$ $SU(2)$, qubit 3 $\leftrightarrow$ $U(1)$) is fixed by the Szangolies construction. Here $|\Omega\rangle$ is a minimal-state ansatz, not a continuum QFT vacuum.

We use $\log(\alpha_i^{-1})$ because it compresses the large dynamic range among inverse couplings and converts multiplicative errors into additive errors. This is a simple first ansatz, not a unique consequence of RG theory. We fixed this mapping before examining any data and did not consider alternatives.

\textbf{Remark (index ambiguity).} In the Szangolies construction the three tensor factors are associated with gauge group structure; Furey's construction associates related algebraic data with generations. No existing framework makes the same three factors carry both meanings simultaneously, so the qubit-gauge assignment above is a modeling choice, not a derived fact. Section~\ref{sec:obstruction} sharpens this: gauge transformations need not act as local unitaries of any chosen qubit factorization, so the factorization itself is gauge-frame dependent. Gauge group, matter representations, generations, Yukawas, and mixing are separate physical targets; the division-algebraic program has made progress on the first and partially the second, and we test only a coupling-hierarchy ansatz here.

\subsection{The Underconstraint Lemma}
\label{sec:underconstraint}

We now prove that the conjecture as stated is underconstrained: unconstrained three-qubit states can reproduce the observed hierarchy without physically selecting a candidate state. This result is the conceptual core of the paper.

\textbf{Lemma 1 (Underconstraint).}\phantomsection\label{lem:under} Under the two-parameter logarithmic map (Eq.~\ref{eq:logmap}), fitting $A$ and $B$ using $\alpha_3$ and $\alpha_2$ reduces the prediction of $\alpha_1$ to a single scalar condition on the entropy triple:
\begin{equation}
r_S \equiv \frac{S_1 - S_2}{S_2 - S_3} = \rSM
\label{eq:ratio}
\end{equation}
where:
\begin{equation}
\rSM = \frac{\log(\alpha_3^{-1}) - \log(\alpha_2^{-1})}{\log(\alpha_2^{-1}) - \log(\alpha_1^{-1})}
     = \frac{2.138 - 3.389}{3.389 - 4.078}
     = 1.817
\label{eq:rsm}
\end{equation}

\textbf{Proof.} From Eq.~(\ref{eq:logmap}), $A + B S_1 = \log(\alpha_3^{-1})$ and $A + B S_2 = \log(\alpha_2^{-1})$. Subtracting: $B(S_1 - S_2) = \log(\alpha_3^{-1}) - \log(\alpha_2^{-1})$. The prediction $A + B S_3 = \log(\alpha_1^{-1})$ requires $B(S_2 - S_3) = \log(\alpha_2^{-1}) - \log(\alpha_1^{-1})$. Dividing these two equations yields Eq.~(\ref{eq:ratio}). \hfill$\square$

The condition $r_S = 1.817$ is a single constraint on the two-dimensional space of ordered entropy triples ($S_1 > S_2 > S_3$) in $[0,1]^3$. This defines a one-parameter family (a curve) in entropy space. (The finer-resolution value $\rSM = 1.8174$ is used as the matching target throughout the computations.)

\subsection{Marginal-Polytope Feasibility}
\label{sec:feasibility}

Each $S_i \in [0, 1]$, and the triangle inequality $S_i \le S_j + S_k$ holds for all tripartite pure states \cite{nielsen}. The sharp feasibility constraint is the polygon inequality of Higuchi, Sudbery, and Szulc \cite{higuchi}: the smaller eigenvalues $\lambda_i$ of the single-qubit reduced density matrices (obtained by solving $S_i = h_2(\lambda_i)$ with $\lambda_i \in [0, 1/2]$) are compatible with a pure three-qubit state if and only if $\lambda_i \le \sum_{j \ne i} \lambda_j$. The required entropy ordering is $S_1 > S_2 > S_3$ (stronger coupling $\leftrightarrow$ more entanglement), with $B < 0$.

\textbf{Proposition 1 (Feasibility and underconstraint).} The observed coupling hierarchy at $M_Z$ is compatible with the one-qubit marginal constraints for pure three-qubit states. Consequently, the unconstrained logarithmic coupling-entropy ansatz is underconstrained and cannot by itself provide evidence for any coupling-entanglement correspondence.

\emph{Example.} The entropy triple $S_1 = 0.875$, $S_2 = 0.730$, $S_3 = 0.650$ gives $r_S = 1.8125$, within the stated tolerance of the gap ratio determined by the observed couplings. Solving $S_i = h_2(\lambda_i)$ gives approximately $\lambda_1 = 0.295$, $\lambda_2 = 0.204$, $\lambda_3 = 0.167$, which satisfy the polygon inequalities with strict inequality for the largest eigenvalue: $0.295 < 0.204 + 0.167 = 0.371$. Therefore a pure three-qubit state with these one-qubit reduced spectra exists. Because the inequalities are strict and $h_2$ is continuous and strictly increasing on $[0, 1/2]$, nearby eigenvalue triples remain feasible, and since $r_S$ varies continuously and nontrivially in this neighborhood, the level set $r_S = \rSM$ locally contains a continuous one-dimensional family of compatible entropy triples. \hfill$\square$

\textbf{Consequence.} The observed coupling hierarchy lies inside the image of the one-qubit marginal polytope under the binary entropy map, and the matching manifold is a continuous family within this image. Finding a matching state in a numerical search is essentially guaranteed and does not constitute evidence for the proposed correspondence. Moreover, even fixing the one-qubit reduced spectra does not generally determine the three-qubit state up to local unitaries: for generic interior points of the marginal polytope, additional local-unitary-inequivalent states share the same spectra \cite{sawicki}. The entropy triple is a very coarse invariant, further reinforcing that coupling matching alone cannot identify a physical candidate state.

The unconstrained existence version of the conjecture is falsified only if the target entropy ratio lies outside the feasible region --- which it does not. Therefore, the meaningful test must be strengthened: after imposing independently motivated algebraic constraints on $|\Omega\rangle$, the conjecture is falsified if no algebraically allowed state produces the observed coupling hierarchy within stated tolerance. Sections~\ref{sec:generic}--\ref{sec:wclass} execute that test.

\subsection{Minimality of Three Qubits}
\label{sec:minimality}

For any \emph{two}-qubit pure state, the Schmidt decomposition forces the two single-qubit reduced density matrices to share the same spectrum, so $S_1 = S_2$ identically. (We verified this numerically as a code check: the maximum of $|S_1 - S_2|$ over $10^5$ random two-qubit pure states is $4 \times 10^{-15}$, i.e., machine precision.) Consequently, a two-coupling version of the logarithmic map with $B \ne 0$ is \emph{impossible}, not merely underconstrained: no two-qubit pure state can encode two unequal couplings under any map of the form of Eq.~(\ref{eq:logmap}).

Three qubits is therefore the smallest system on which the coupling-entropy ansatz has any content at all. This is a modest but genuine strengthening of the division-algebra motivation: the octonionic (three-qubit) level of the Hopf hierarchy is not merely the level at which the Standard Model gauge group appears --- it is also the first level at which a coupling hierarchy is even expressible in single-qubit entropies.

\subsection{What Would Constitute Evidence}
\label{sec:evidence}

The underconstraint result establishes that coupling-entropy matching alone is not evidential. Evidence for a physical coupling-entanglement correspondence would require one or more of:

\textbf{Requirement A (Algebraic selection).} The candidate state $|\Omega\rangle$ must not be arbitrary: it must be selected by the Szangolies construction, by an embedding of Furey's Clifford-algebraic structures (constructed in Section~\ref{sec:jw}), or by another independently motivated algebraic condition. If such a constrained state automatically produced the correct entropy gap ratio (Eq.~\ref{eq:ratio}), without fitting, that would be non-trivial. (Section~\ref{sec:octonionic} reports that no canonical construction considered here satisfies this; Appendix~\ref{sec:rotor} reports a dynamical candidate, which is symmetric instead.)

\textbf{Requirement B (Holdout predictions).} A state selected by coupling matching must predict independent observables not used in fitting: the Weinberg angle, the GUT-scale convergence energy, qualitative CKM/PMNS structure, a specific entanglement class, a notable 3-tangle value, or non-generic local-unitary invariants.

\textbf{Requirement C (Non-genericity).} The coupling-matching states must occupy a distinguished, measure-zero or algebraically special subset of state space. (Section~\ref{sec:exp1} reports that this fails empirically for unconstrained matching.)

\section{A Sampling Sanity Check: Haar Genericity and Null Controls}
\label{sec:generic}

The analytic argument of Section~\ref{sec:static} is the primary result of the static analysis: one scalar condition imposed inside a feasible region of entropy triples yields a positive-dimensional matching set, so matches exist and carry no evidence. This section reports a numerical confirmation --- deliberately framed as a supporting sanity check, not a discovery --- that the matching manifold is also statistically unremarkable, together with four null controls. (Every numerical claim in this paper was independently re-verified in a single reproducibility run, including a cross-check of the coupling inputs against PDG primary values $\alpha_{\mathrm{em}}^{-1} = 127.951$, $\sin^2\theta_W = 0.23122$, which reproduces $\alpha_1$ to within 0.04\% and $\alpha_2$ to within 0.2\%, the residuals reflecting scheme conventions.) Summary plots appear in Figure~\ref{fig:experiments}.

\begin{figure}[t]
\centering
\includegraphics[width=\textwidth]{fig1_experiments.pdf}
\caption{Summary of the computational program. (a) Scale dependence of the gap ratio $\rSM(\mu)$ under one-loop Standard Model running, which varies by $\sim$24\% between $M_Z$ and $10^6$ GeV and becomes singular near $10^{13}$ GeV, where $\alpha_1^{-1}$ and $\alpha_2^{-1}$ cross (Section~\ref{sec:running}). (b) The matching manifold is statistically generic: the 3-tangle distribution of coupling-matched states (from $10^7$ Haar-random states) is nearly identical to that of ordering-conditioned Haar controls (Kolmogorov-Smirnov $D = 0.031$; Section~\ref{sec:exp1}). (c) The weighted-W matching curve: a continuous one-parameter family of W-class ($\tau_3 = 0$) states satisfying $r_S = 1.8174$ exactly --- given in closed form by Lemma~2 (Section~\ref{sec:wfamily}) --- showing that coupling matching does not fix the SLOCC class (Section~\ref{sec:wfamily}).}
\label{fig:experiments}
\end{figure}

\subsection{Haar Genericity (Experiment 1)}
\label{sec:exp1}

Among $N = 10^7$ Haar-random three-qubit states (labels preserved per the Szangolies assignment), 1{,}666{,}906 had the required ordering $S_1 > S_2 > S_3$ ($\approx 1/6$, as exchange symmetry requires) and 2{,}666 satisfied $|r_S - 1.8174| < 0.01$ --- a match rate of $2.7 \times 10^{-4}$ overall. Matching is easy, exactly as Lemma~1 predicts. The matched states are statistically generic: their 3-tangle \cite{coffman} distribution is nearly identical to that of ordering-conditioned Haar controls (Kolmogorov-Smirnov $D = 0.031$, a negligible effect size despite $p = 0.023$ at these sample sizes; medians 0.326 versus 0.312; Figure~\ref{fig:experiments}b), and the fraction of states with $\tau_3 < 0.01$ is approximately 0.1\% in both populations, so the matching condition does not concentrate states near the W-class locus or any other special set. (Pairwise concurrence distributions do shift, KS up to 0.31, but only as forced by conditioning on entropy gaps; higher local-unitary invariants are left to future work.) \textbf{Requirement C (non-genericity, Section~\ref{sec:evidence}) therefore fails empirically as well as analytically.}

\subsection{Null Controls}
\label{sec:controls}

Four controls confirm that parameter freedom, not algebraic structure, drives matching.

\textbf{Control A (Random coupling targets).} Matched states for the measured target have mean $\tau_3 = 0.320$ versus $0.337 \pm 0.018$ for randomly sampled coupling triples (log-uniform on $[0.001, 0.5]$, ordered) --- a deviation of $z = -0.92$, within statistical noise: the measured hierarchy is indistinguishable from a random target as far as its matching manifold is concerned.

\textbf{Control B (Two qubits).} The two-qubit map is not merely weaker --- it is impossible: $S_1 = S_2$ identically (Section~\ref{sec:minimality}, which emerged from this control).

\textbf{Control C (Four qubits).} Four-qubit states, which correspond to no normed division algebra, match at least as easily: per-triple rate $3.4 \times 10^{-4}$ versus $2.6 \times 10^{-4}$ for three qubits in the same $N = 10^6$ control run, and $1.3 \times 10^{-3}$ if any of the four ordered triples may match.

\textbf{Control D (Shuffled labels).} All six qubit-gauge permutations match at statistically identical rates (256--292 matches per $10^6$ states), as the exchange symmetry of the Haar measure requires; the Szangolies labeling carries no measurable content for unconstrained states. Section~\ref{sec:obstruction} explains why: the labeling is not gauge-invariant in the first place.

\section{The Weighted-W Counterexample and the Rigidity of Canonical Octonionic States}
\label{sec:wclass}

\subsection{An Analytic Matching Family in the W Class}
\label{sec:wfamily}

The deepest failure of Requirement C is not statistical but structural: the matching manifold intersects the W class, the SLOCC class with identically zero 3-tangle. The computational program found this as a numerical curve (Figure~\ref{fig:experiments}c); here we upgrade it to a closed-form result, so that the counterexample no longer rests on numerics.

For a weighted-W state $|\psi\rangle = a|001\rangle + b|010\rangle + c|100\rangle$ (real $a, b, c \ge 0$, $a^2 + b^2 + c^2 = 1$), the single-qubit reduced density matrices are exactly diagonal: $\rho_1 = \mathrm{diag}(a^2 + b^2,\, c^2)$, $\rho_2 = \mathrm{diag}(a^2 + c^2,\, b^2)$, $\rho_3 = \mathrm{diag}(b^2 + c^2,\, a^2)$. Hence
\[
S_1 = h_2(c^2), \qquad S_2 = h_2(b^2), \qquad S_3 = h_2(a^2),
\]
with $h_2$ the binary entropy, and $\tau_3 = 0$ identically on the family. The entropies are therefore computable in closed form from the weights, and the matching condition becomes a scalar equation in the weights alone.

\textbf{Lemma 2 (Analytic weighted-W matching family).}\phantomsection\label{lem:wfamily} Let $r > 1$ be a target gap ratio and let $h_2^{-1} : [0,1] \to [0, 1/2]$ denote the inverse of the binary entropy on $[0, 1/2]$. Let $g_{\max}(r)$ be the unique root of
\[
h_2^{-1}(1 - r g) + h_2^{-1}(1 - (1+r) g) = \tfrac{1}{2}.
\]
For each $g \in (0, g_{\max}(r))$ there exists a unique $s = s(g) \in (g,\, 1 - rg)$ solving
\begin{equation}
h_2^{-1}\big(s + r g\big) + h_2^{-1}(s) + h_2^{-1}\big(s - g\big) = 1,
\label{eq:wfamily}
\end{equation}
and the weighted-W states $|\psi(g)\rangle$ with
\[
c^2 = h_2^{-1}\big(s(g) + rg\big), \qquad b^2 = h_2^{-1}\big(s(g)\big), \qquad a^2 = h_2^{-1}\big(s(g) - g\big)
\]
form a smooth, injective one-parameter curve of W-class states ($\tau_3 \equiv 0$) with $S_1 = s + rg > S_2 = s > S_3 = s - g$ and
\[
r_S = \frac{(s + rg) - s}{s - (s - g)} = r \quad \text{exactly, for every } g.
\]
As $g \to 0^+$ the curve terminates at the W state ($s \to h_2(1/3) \approx 0.9183$); as $g \to g_{\max}^-$ it terminates at the maximal-asymmetry endpoint with $S_1 = 1$ ($c^2 = 1/2$). For the measured target $r = \rSM = 1.8174$: $g_{\max} = 0.082561$, and the endpoint weights are $(a^2, b^2, c^2) = (0.2240, 0.2760, 0.5000)$.

\textbf{Proof.} Existence and uniqueness: define $\Phi(s, g)$ as the left side of Eq.~(\ref{eq:wfamily}) minus 1. Since $h_2^{-1}$ is strictly increasing, $\Phi$ is strictly increasing in $s$ on $(g, 1 - rg)$. At the lower endpoint, $\Phi(g^+, g) = h_2^{-1}((1+r)g) + h_2^{-1}(g) - 1 < 0$ because each term is strictly less than $1/2$. At the upper endpoint, $\Phi((1 - rg)^-, g) = \tfrac12 + h_2^{-1}(1 - rg) + h_2^{-1}(1 - (1+r)g) - 1$, which is strictly positive precisely when $g < g_{\max}(r)$, since both terms are strictly decreasing in $g$ and vanish jointly past the defining root. The intermediate value theorem and strict monotonicity give a unique root $s(g)$; the implicit function theorem ($\partial\Phi/\partial s > 0$, continuously differentiable in the interior) gives smoothness of $g \mapsto s(g)$, hence of the curve. The state assignment realizes the prescribed marginals because Eq.~(\ref{eq:wfamily}) is exactly the normalization $a^2 + b^2 + c^2 = 1$ and the marginals are diagonal as computed above; $S_1 > S_2 > S_3$ holds by construction, and $r_S = rg/g = r$. Injectivity: $g \mapsto S_2 - S_3 = g$ is injective. \hfill$\square$

\textbf{Numerical verification.} We verified Lemma~2 against the partial-trace machinery of the computational program. For $g \in \{0.005, 0.02, 0.04, 0.06, 0.0779, 0.082\}$ the reconstructed states give $r_S = 1.8174000000$ to ten digits, and the curve reproduces the five stored matching-curve grid solutions of the original experiment (e.g.\ $S = (0.9995, 0.8580, 0.7802)$ at $g = 0.0779$) to within the original grid resolution ($\le 3.4 \times 10^{-6}$ in the weights). The verification script is \path{experiments/verify_wfamily.py}; the stored grid solutions are in \path{experiments/results/exp2_results.json}.

\textbf{Consequence.} Coupling matching does not select GHZ-type tripartite entanglement, is not tied to the Cayley hyperdeterminant (which vanishes identically on the family), and is not tied to any special octonionic locus: an explicit closed-form W-class family satisfies the ratio condition exactly. Combined with Section~\ref{sec:exp1}, the static conjecture fails Requirement C both statistically and structurally. (A follow-up computation shows this family is, in addition, pure gauge under the gauge-invariant accounting --- see the Note Added.)

\subsection{Canonical Octonionic Constructions Are Permutation-Degenerate}
\label{sec:octonionic}

The complementary question is whether any independently motivated algebraic state lands on the matching manifold (Requirement A). We tested the natural candidates: the GHZ and W reference families; the weighted-W family above; octonion-structure states (amplitudes from Fano-plane incidence data and multiplication-table structure constants, under several inequivalent maps); quaternionic-line states for all seven Fano lines; and the preferred-complex-direction states $(e_0 + i e_\ell)/\sqrt{2}$ of the Szangolies dimensional reduction. The octonion multiplication table was generated by Cayley--Dickson doubling and verified numerically to satisfy $e_a e_b = \pm e_{a \oplus b}$; convention-independence was checked separately (Appendix~\ref{sec:robustness}, and an independent reconstruction from a Fano-triple table recorded in the repository's verification log).

\begin{center}
\small
\begin{tabular}{p{4.5cm}p{4.3cm}cp{2.6cm}}
\hline
Family & Entropies $(S_1, S_2, S_3)$ & $\tau_3$ & $r_S = 1.8174$ reachable? \\
\hline
GHZ($\theta$) & $S_1 = S_2 = S_3$ & $4\sin^2\theta\cos^2\theta$ & No ($r_S$ undefined) \\
W & $(0.918, 0.918, 0.918)$ & 0 & No ($r_S$ undefined) \\
Weighted W & tunable & 0 & \textbf{Yes --- closed-form curve (Lemma~2)} \\
Octonion sign-row / sign-convolution map & $(0.811, 0.811, 0.811)$ & 0.25 & No \\
Octonion diagonal map $\sigma(a,a)$ & $(0.811, 0.811, 0.811)$ & 0.25 & No \\
Fano incidence (uniform over $e_1 \ldots e_7$) & $(0.349, 0.349, 0.349)$ & 0.082 & No \\
Uniform over $e_0 \ldots e_7$ & $(0, 0, 0)$, separable & 0 & No \\
Quaternionic-line states (all 7 Fano lines) & degenerate: $(0,0,0)$, $(1,1,0)$ perms, or $(1,1,1)$ & 0 or 1 & No \\
Preferred-complex-direction $(e_0 + i e_\ell)/\sqrt{2}$, $\ell = 1 \ldots 7$ & degenerate: $(0,0,0)$, $(1,1,0)$ perms, or $(1,1,1)$ & 0 or 1 & No \\
\hline
\end{tabular}
\end{center}

\textbf{Every canonical construction considered here yields permutation-degenerate entropies} --- all three equal, or degenerate $\{0,0,0\}$/$\{1,1,0\}$/$\{1,1,1\}$ patterns. None can produce a three-way hierarchy; $r_S$ is undefined ($0/0$) on all of them. The algebra's symmetry ($G_2$ automorphisms, triality, the XOR-grading of the basis) is too rigid: producing a hierarchy requires symmetry breaking beyond the choice of a preferred unit. This is a scoped statement --- we enumerate the constructions tested in the table, and make no claim about constructions outside it. It sharpens Requirement A: whatever selection principle the framework eventually supplies, it cannot be any of the obvious ``canonical state'' maps from octonionic structure data. (Appendix~\ref{sec:rotor} explores the dynamical route and finds the same rigidity.)

Note: Furey's Clifford-algebraic construction \cite{furey14, furey18} does not directly produce three-qubit states. At the time the program was designed, an embedding from $\mathrm{Cl}(6)$ structures to $\CC^2 \otimes \CC^2 \otimes \CC^2$ appeared to be an open problem. It is not: the Jordan--Wigner transformation provides the canonical isomorphism, and Section~\ref{sec:jw} develops its consequences.

\section{The Jordan--Wigner Gauge-Covariance Obstruction}
\label{sec:jw}

\subsection{The Embedding Is Standard Mathematics Made Explicit}
\label{sec:embedding}

Furey's construction is explicitly built from three fermionic ladder operators \cite{furey14, furey18}, and the minimal-left-ideal structure of $\mathrm{Cl}(6)$ with electric charge as one-third the number operator is developed in detail by Stoica \cite{stoica}. The complex Clifford algebra $\mathrm{Cl}(6)$ is isomorphic to $M_8(\CC)$, which is exactly the full operator algebra of three qubits, and the Jordan--Wigner transformation \cite{jordanwigner} between three fermionic modes and three qubits is textbook. The accompanying code verifies that the Jordan--Wigner images of the three ladder operators satisfy the canonical anticommutation relations exactly (Experiment 5; \texttt{experiments/results/exp5\_results.json}).

The bridge between the Furey and Szangolies constructions is therefore a standard transformation, dating to 1928, that had not been put to work in this specific context. What this paper contributes is the reading of its consequences --- and, principally, the obstruction it exposes (Section~\ref{sec:obstruction}).

\subsection{Computed Consequences}
\label{sec:jwconsequences}

\begin{itemize}

\item \textbf{Charge is Hamming weight.} Furey's minimal-left-ideal basis for one generation --- $\nu$, three $\bar{d}$ colors, three $u$ colors, $e^+$ --- maps under Jordan--Wigner to the three-qubit computational basis, graded by Hamming weight: $\nu$ at weight 0, the $\bar{d}$ triplet at weight 1, the $u$ triplet at weight 2, $e^+$ at weight 3, with electric charge $Q = (\text{weight})/3$. Furey's observation that charge quantization follows from the integer spectrum of a number operator \cite{furey18} becomes, in qubit language, the statement that charge quantization is the integer grading of the computational basis.

\item \textbf{SLOCC classes acquire particle-physics meaning.} The color-symmetric quark combination (uniform superposition over the weight-1 triplet) is precisely the \textbf{W state}, with $S = (0.918, 0.918, 0.918)$ and $\tau_3 = 0$. The $\nu + e^+$ superposition is precisely the \textbf{GHZ state} --- the maximal violation of charge superselection ($\Delta Q = 1$) maps to maximal genuine tripartite entanglement.

\item Fock (particle) basis states are computational basis states: single particles carry no entanglement in this frame.

\end{itemize}

\subsection{Critical Caveat: Entanglement Is Gauge-Frame Dependent}
\label{sec:obstruction}

The same computation exposes a structural obstruction. Color rotations act on the fermionic \emph{modes}, not on the qubit tensor factors. A discrete Fourier transform in color space --- an $SU(3)$ gauge transformation --- maps a separable single-quark state, $S = (0, 0, 0)$, to a state with $S = (0.918, 0.918, 0.918)$. A numerical check confirms that this transformation is not a local unitary of the qubit factorization.

That entanglement is relative to a choice of tensor-product structure, and that mode rotations change it, is well documented in the quantum-information literature \cite{zanardi, barnum}; we are not claiming the phenomenon as new. The contribution here is its specific consequence for this program:

\textbf{Gauge transformations are not local unitaries of the qubit factorization, and therefore single-qubit entanglement entropies in this frame are gauge-frame dependent.} Any physical coupling-entropy correspondence must consequently be built from gauge-invariant entanglement functionals (notion (ii) of Table~\ref{tab:notions}), or must include an algebraic gauge-fixing of the tensor factorization as part of its definition. Neither has yet been constructed within the static program. Until one is, the coupling-entropy question is ill-posed as stated --- which is itself the sharpest result the executed program produced. In one line: \emph{the naive entropy observable is not gauge-invariant; therefore the map is not physically meaningful.}

This obstruction sharpens, and partially explains, two earlier findings: the index ambiguity (Section~\ref{sec:conjecture}, Remark) --- the tensor-factor labels were never gauge-covariantly defined --- and the Control D null (Section~\ref{sec:controls}): no unconstrained entanglement statistic could have privileged one qubit-gauge labeling, because the labeling is not gauge-invariant in the first place. We regard the construction of the right gauge-invariant entanglement functional as the central open problem this paper leaves for the coupling-entropy program. A follow-up computation carries out exactly this gauge-invariantization for the minimal toy and sharpens the present conclusion --- see the Note Added preceding the references.

\section{The Renormalization-Group Obstruction and the Entropy-Flow Consistency Envelope}
\label{sec:rg}

The obstruction of Section~\ref{sec:jw} is the principal reason the static map fails: its observable is not gauge-invariant. This section records a second, logically independent inconsistency --- with renormalization-group running --- and distills what survives of a flow version of the map into a consistency envelope. The material is deliberately brief. The inconsistency itself is true but not deep, and the envelope constrains only map parameters, not observables.

\subsection{Scale Dependence of the Gap Ratio}
\label{sec:running}

The one-loop RG equations, with the convention $dg_i/d(\log\mu) = b_i g_i^3/(16\pi^2)$, give:
\begin{equation}
\alpha_i^{-1}(\mu) = \alpha_i^{-1}(M_Z) - \frac{b_i}{2\pi}\,\log\frac{\mu}{M_Z}
\label{eq:running}
\end{equation}
\[
(b_1, b_2, b_3) = \left(\tfrac{41}{10},\, -\tfrac{19}{6},\, -7\right)
\]

The measured couplings run, and the gap ratio runs with them: $\rSM(\mu)$ falls from 2.072 at 10 GeV through 1.817 at $M_Z$ to 1.385 at $10^6$ GeV --- a $\sim$24\% variation between $M_Z$ and $10^6$ GeV --- and the denominator crosses zero near $10^{13}$ GeV, where $\alpha_1^{-1}$ and $\alpha_2^{-1}$ cross (Figure~\ref{fig:experiments}a). A fixed entropy triple satisfying $r_S = 1.817$ at $M_Z$ therefore cannot match the coupling hierarchy at any other scale: the logarithmic map with scale-independent entropies is inconsistent with RG running. Either (a) the logarithmic map is the wrong ansatz, or (b) the entropies themselves must run, $S_i = S_i(\mu)$, or (c) the correspondence holds only at one privileged scale --- and the divergence near $10^{13}$ GeV is not unification ($\alpha_1$ and $\alpha_2$ cross while $\alpha_3$ remains distinct; the non-supersymmetric model does not unify at one loop). The next subsection records the sharpest defensible form of option (b).

\subsection{The Entropy-Flow Consistency Envelope}
\label{sec:envelope}

\textbf{This subsection is a consistency envelope, not a surviving conjecture.} The affine map below does not derive the entropy flow; it relocates the question from ``why this entropy triple?'' to ``why this entropy trajectory?''. What it delivers is the set of constraints that \emph{any} viable flow version of the correspondence would have to satisfy. No model satisfying them is proposed here.

\textbf{Affine flow map.} There exist real constants $A$ and $B$ (with $B < 0$) and a one-parameter family of three-qubit pure candidate states $|\Omega(\mu)\rangle$ such that
\begin{equation}
\alpha_i^{-1}(\mu) = A + B \cdot S_i(\mu) \qquad \text{for } i = 1, 2, 3
\label{eq:flowmap}
\end{equation}
with the qubit-gauge assignment as before. Substituting the one-loop running (Eq.~\ref{eq:running}) gives
\begin{equation}
\frac{dS_i}{d\log\mu} = -\frac{b_i}{2\pi B}
\label{eq:flow}
\end{equation}
so the entropies run \emph{linearly} in $\log\mu$, with slopes proportional to the one-loop beta coefficients. The form is RG-consistent by construction, and the slopes admit a counting reading --- beta coefficients count charged degrees of freedom, and entanglement entropy in critical systems runs with coefficients that count degrees of freedom (Zamolodchikov \cite{zamolodchikov}; Casini--Huerta \cite{casini}; cf.\ entanglement renormalization \cite{vidal}) --- but this reading remains entirely open.

\textbf{Honesty clause.} The flow map is a reparametrization of measured couplings: given any running $\alpha_i^{-1}(\mu)$ and any $(A, B)$, Eq.~(\ref{eq:flowmap}) \emph{defines} a trajectory $S_i(\mu)$. Its only possible content lives in (i) the quantum-marginal feasibility constraints on $(A, B)$ below --- which constrain only the unobservable map parameters; (ii) the boundary condition at the symmetric point (Section~\ref{sec:corollary}), its single falsifiable consequence; and (iii) the counting interpretation of the slopes, where any genuine physical content would have to reside. We do not claim uniqueness of the affine form (per-coupling constants $(A_i, B_i)$ are equally RG-consistent and even less constrained).

\textbf{Computed feasibility constraints} (Experiment 3; results in \path{experiments/results/exp3_results.json}). For the trajectory $S_i(\mu)$ to be realizable by pure three-qubit states at every $\mu \in [M_Z, M_{\mathrm{Planck}}]$, the constants $(A, B)$ must satisfy $S_i(\mu) \in [0, 1]$ together with the Higuchi--Sudbery--Szulc polygon inequality \cite{higuchi} at every scale. Scanning the $(A, B)$ plane:

\begin{itemize}

\item The box constraint $S \in [0,1]^3$ alone requires $|B| \gtrsim 51$. Adding the polygon inequality raises this to $\boldsymbol{|B| \gtrsim 126.5}$ --- robust against enlarging the scan region and against moving the UV endpoint from $M_{\mathrm{Planck}}$ down to $10^{16}$ GeV. The feasible region itself is unbounded (larger $A$ forces proportionally larger $|B|$); the lower bound on $|B|$ is the entire content of the constraint.

\item Two honest deflations, identified in independent verification. First, the constraints \emph{bind at $M_Z$ alone}: imposing feasibility only at $M_Z$ yields the identical feasible $(A, B)$ set, because the coupling spread is widest there. The flow adds no restrictive power beyond the static affine map at one scale. Second, since the feasible set is non-empty, nothing observable is constrained.

\item At the minimal-$|B|$ feasible point, the three entropies are confined to a mid-band ($\approx 0.47$--$0.87$) at all scales, with small slopes $|dS_i/d\log\mu| \approx 0.004$--$0.009$ per e-fold.

\item Structural feature (a restatement, not a prediction): the entropy ordering inverts in the UV. At $M_Z$, $S_1 > S_2 > S_3$; above the pairwise coupling crossings --- $9.7 \times 10^{12}$ GeV ($\alpha_2$--$\alpha_1$), $2.4 \times 10^{14}$ GeV ($\alpha_3$--$\alpha_1$), $1.05 \times 10^{17}$ GeV ($\alpha_3$--$\alpha_2$) --- the ordering reverses and the $U(1)$ factor becomes the most entangled. This is the known one-loop crossing pattern re-expressed in entropy variables.

\end{itemize}

These results are summarized in Figure~\ref{fig:flow}.

\begin{figure}[t]
\centering
\includegraphics[width=\textwidth]{fig2_entropy_flow.pdf}
\caption{The entropy-flow consistency envelope (Section~\ref{sec:envelope}). (a) Running entropies $S_i(\mu)$ under Eq.~(\ref{eq:flowmap}) at a representative feasible point $(A, B)$: linear in $\log\mu$ with slopes proportional to the one-loop beta coefficients; the ordering inverts above the pairwise coupling crossings. (b) The feasible $(A, B)$ region: box constraint alone (light) versus box plus Higuchi-Sudbery-Szulc polygon inequality (dark), which raises the bound to $|B| \gtrsim 126.5$; the star marks the representative minimal-$|B|$ point. (c) Entanglement asymmetry $\Delta S(\mu)$: the Standard Model never symmetrizes --- $\Delta S$ bottoms out at $\approx 0.029$ near $2.4 \times 10^{14}$ GeV and rises again --- while the MSSM (1 TeV threshold) cusps to $\approx 0.003$ near $1.4 \times 10^{16}$ GeV.}
\label{fig:flow}
\end{figure}

\subsection{The Unification Corollary and the Normalization Caveat}
\label{sec:corollary}

Under the flow map, exact entanglement symmetry $S_1 = S_2 = S_3$ at some scale is \emph{equivalent} to exact coupling unification at that scale --- immediately, from Eq.~(\ref{eq:flowmap}) with $B \ne 0$. The equivalence is definitional. Define the entanglement asymmetry $\Delta S(\mu)$ as the spread of the inverse couplings divided by $|B|$. Computed at one loop:

\begin{itemize}
\item \textbf{Standard Model:} $\Delta S(\mu)$ is minimized at $\mu^* = 2.4 \times 10^{14}$ GeV with $\Delta S \approx 0.029$ (at the minimal feasible $|B|$) --- and it \emph{never reaches zero} (Figure~\ref{fig:flow}c). Note also that $\mu^*$ is an \emph{interior} scale: above it the couplings re-split, so a symmetric boundary condition would sit in the middle of the flow, not at its UV end.
\item \textbf{MSSM (1 TeV threshold):} $\Delta S$ cusps to $\approx 0.003$ at $1.4 \times 10^{16}$ GeV, reflecting the well-known approximate unification of MSSM couplings \cite{amaldi}.
\end{itemize}

The corollary, stated conditionally: \textbf{if} one demands an exactly symmetric boundary state, \textbf{then} the map requires exact coupling unification, which the Standard Model alone does not provide. This is the envelope's single falsifiable consequence, and it is inherited rather than novel: it restates the unification question.

\textbf{The normalization caveat, again, because it bites hardest here.} All of these numbers use $\alpha_1 = (5/3)\alpha_Y$ (Section~\ref{sec:conjecture}). The factor $5/3$ is motivated by GUT embedding; a framework that obtains $SU(3) \times SU(2) \times U(1)$ directly from octonionic structure --- without a GUT group --- has no derivation of it, and without it, three trajectories meeting at a point would be an unexplained codimension-one coincidence. The unification corollary is therefore doubly conditional: on the symmetric-boundary reading, and on a normalization that only a GUT-type embedding currently motivates. (Possible internal routes to fixing the normalization --- anomaly cancellation, trace normalization over representation content, division-algebraic inner products --- are listed as open problems; none is developed here.)

\textbf{A point of contact with Appendix~\ref{sec:rotor}.} If the symmetric boundary state is taken to be the rotor vacuum of Appendix~\ref{sec:rotor} (entropies $h_2(1/7) \approx 0.5917$), Eq.~(\ref{eq:flowmap}) at the symmetric scale fixes one linear relation, $A + 0.5917\,B = \alpha_*^{-1}$, and feasible $(A, B)$ satisfying it exist for every unification coupling $\alpha_*^{-1}$ in the plausible range $[20, 50]$. The rotor boundary state is therefore \emph{compatible} with the envelope but does not constrain it --- consistent with the recurring pattern that internal consistency is cheap and constraint is what must be earned.

\section{Discussion}
\label{sec:discussion}

\subsection{What Is Ruled Out and What Is Not}
\label{sec:notruledout}

All no-go results above are scoped (Section~\ref{sec:scope}). Table~\ref{tab:notruledout} states the scope explicitly.

\begin{table}[ht]
\centering
\begin{tabular}{p{6.6cm}|p{6.6cm}}
\hline
\textbf{Ruled out (in this paper's scope)} & \textbf{Not ruled out} \\
\hline
Static single-qubit entropy matching of the measured couplings (underconstrained, generic, scale-inconsistent) & Gauge-invariant flux-sector or charge-sector statistics as coupling carriers \\
\hline
Unconstrained three-qubit matching as evidence (matching is guaranteed and SLOCC-blind; Lemma~2) & Algebra-selected states with independent, holdout predictions (Requirements A--B) \\
\hline
Frame-dependent Jordan--Wigner single-qubit entropy as a physical observable & Algebra-relative (gauge-invariant) entanglement functionals \\
\hline
A scale-independent entropy triple under the logarithmic map & An RG-consistent entropy/flux flow satisfying the envelope of Section~\ref{sec:envelope} \\
\hline
Hierarchy from the canonical octonionic constructions considered here (kinematic and dynamical: Section~\ref{sec:octonionic} and Appendix~\ref{sec:rotor}) & Hierarchy from symmetry-breaking structure beyond those constructions \\
\hline
\end{tabular}
\caption{Scope of the no-go results. The left column is established by this paper, within the stated ansatz and the enumerated state and Hamiltonian families. The right column is untouched by it.}
\label{tab:notruledout}
\end{table}

\subsection{Relation to Subsequent Work}
\label{sec:series}

This paper is self-contained, and nothing in it depends on unpublished work; we confine the pointers to essentials. The gauge-invariant repair that Section~\ref{sec:jw} calls for has since been carried out for the minimal Jordan--Wigner toy: entanglement relative to the gauge-invariant operator subalgebras (notion (ii) of Table~\ref{tab:notions}) can be computed exactly there, and the outcome --- including the fate of the weighted-W matching curve of Lemma~2 --- is summarized in the Note Added preceding the references. Separately, in one exactly solvable setting (two-dimensional Yang--Mills), a gauge coupling \emph{is} genuinely reconstructible from gauge-invariant edge/flux-sector statistics (notion (iii)); that positive result, and the sharp distinction between \emph{reconstruction} within a presupposed theory and \emph{determination} of a coupling by entanglement, is developed in a separate manuscript. For the avoidance of doubt: no claim that real four-dimensional couplings are determined by entanglement is made here or in that work; that proposition remains unsupported speculation.

\subsection{Broader Motivations, Set Aside}
\label{sec:broader}

Earlier versions of this manuscript reviewed the broader information-first landscape --- de Sitter holography, the cosmological constant, the monopole problem under a GUT-free reading of the gauge group, and emergent time via Page--Wootters \cite{page} --- as context for the coupling-entropy question. None of those topics is materially advanced by the computations reported here, so we record them in this one paragraph and claim nothing about them. The single point of contact that survives contact with the technical results is the normalization issue (Sections~\ref{sec:conjecture} and \ref{sec:corollary}): a framework that wants gauge structure without GUT embedding owes an independent derivation of the abelian normalization, and until one exists, every $\alpha_1$-dependent quantity in this program is partly conventional.

\subsection{Reproducibility}
\label{sec:repro}

\begin{sloppypar}
All computations are reproducible from the repository published with this paper. The \path{experiments/} directory contains: the shared library (\path{qio_lib.py}: Haar sampling via complex Gaussian normalization, partial traces via tensor reshaping, entropies from eigenvalues of reduced density matrices, 3-tangle and concurrences); one script per experiment (\path{exp1_matching_manifold.py}, \path{exp2_algebraic_families.py}, \path{exp3_entropy_flow.py}, \path{exp4_dynamics_toys.py}, \path{exp5_furey_jw.py}, \path{exp6_rotor_robustness.py}), the null controls (\path{controls.py}), the figure generators (\path{make_figures.py}, \path{make_fig2.py}), the verification script for Lemma~2 (\path{verify_wfamily.py}), and a cross-check of every number quoted in the text (\path{verify_numbers.py}). All scripts fix their random seeds, run on NumPy/SciPy only, and write their raw outputs as JSON to \path{experiments/results/}; the exact output files used for this manuscript are committed alongside the code. The repository also includes an independent verification log and a response-to-review document.
\end{sloppypar}

\section{Conclusion}
\label{sec:conclusion}

We asked whether the measured three-gauge-coupling hierarchy is encoded, under a proposed three-qubit entropy map, in the entanglement of a three-qubit candidate state. The answer is no in every form we could make precise, and the results are scoped no-gos, each verified with published code. Analytically: the static map is underconstrained, collapsing to one gap-ratio condition imposed inside a feasible region, so the matching set is positive-dimensional and matching is not evidence (Lemma~1, Proposition~1); two-qubit states cannot encode unequal couplings at all; and an explicit closed-form W-class family satisfies the ratio exactly (Lemma~2), so matching does not even fix the SLOCC class. A sampling check across $10^7$ Haar states confirms this picture: the matching manifold is statistically generic, and four null controls confirm that parameter freedom drives matching. Deepest: gauge transformations are not local unitaries of the qubit factorization, so the quantity the conjecture is about is gauge-frame dependent --- \emph{the naive entropy observable is not gauge-invariant, and therefore the map is not physically meaningful as stated}. Independently, a scale-independent entropy triple is inconsistent with RG running, the flow version surviving only as a consistency envelope that constrains map parameters, not observables. And as a control: every canonical octonionic construction considered here is permutation-degenerate, and the exactly solvable rotor Hamiltonian $i(L_u + 2R_u)$ of Appendix~\ref{sec:rotor} --- the one genuine vacuum in this paper --- selects a convention-invariant, permutation-symmetric ground state with marginal spectra $(1/7, 6/7)$ and 3-tangle $8/49$: even the most canonical dynamical octonionic construction produces symmetry, not hierarchy.

To restate the scope plainly: we are not refuting Szangolies's construction, which concerns gauge-group structure and is untouched by these results; we are refuting a natural but naive coupling-entropy extension of it. What survives is not a model but a specification: any future coupling-entanglement correspondence must be built from gauge-invariant (algebra-relative) entanglement functionals (notion (ii) of Table~\ref{tab:notions}), must generate hierarchy by flow within the consistency envelope of Section~\ref{sec:envelope} (whose only falsifiable consequence restates the unification question and inherits the abelian normalization convention), and must explain why the algebra's own canonical vacuum is symmetric. A map of where the answers cannot lie is what this paper contributes; we offer it to whoever next enters this territory.

\appendix

\section{An Exact Octonionic Rotor Construction as a Dynamical Control}
\label{sec:rotor}

This appendix reports a result of a different kind from the no-go results of the main text, and we frame it explicitly as a control on the whole program: \emph{even the most canonical dynamical octonionic construction produces symmetry, not hierarchy.} Throughout the main text no principle selected the candidate state; here, a Hamiltonian canonically constructed from the octonion algebra --- with no tuned parameters --- selects its unique ground state. \textbf{This is the one place in the paper where ``vacuum'' is meant literally: the state is selected by algebra-canonical dynamics, not by fit.} The verdict of the control is deflationary in the same direction as everything above: the selected vacuum is permutation-symmetric, not hierarchical.

\subsection{The Hamiltonian and Its Closed Form}
\label{sec:closedform}

\textbf{Status: computed (Experiments 4 and 6; results in \path{experiments/results/exp4_results.json} and \path{exp6_results.json}); the closed form makes all reported invariants exact.}

The left-multiplication operators $L_i = L_{e_i}$ of the imaginary octonion units on $\OO \cong \RR^8$ satisfy the Clifford relation $L_i L_j + L_j L_i = -2\delta_{ij}$ (verified numerically in our run), so the operators $iL_i$ are Hermitian on $\CC^8$ --- which is exactly the three-qubit Hilbert space. This permits Hamiltonians built canonically from the algebra.

\textbf{Linear no-go.} Every linear Hamiltonian $H = \sum_i c_i\,(iL_i)$ satisfies $H^2 = |c|^2\,\mathbf{1}$, so its spectrum is $\pm|c|$ with each level 4-fold degenerate, for \emph{every} coupling vector $c$. The Clifford structure forbids linear algebra-canonical dynamics from selecting a unique vacuum; within the 4-dimensional ground manifold, entropy triples are unconstrained by the dynamics.

\textbf{Commuting quadratic pairings.} The 105 ways of partitioning six of the seven generators into three commuting quadratic pairs give discrete ground states whose sorted entropy triples take only four values --- $(0, 0, 0)$, $(0, 1, 1)$, $(0.811, 0.811, 0.811)$, and $(1, 1, 1)$ --- all separable, degenerate, or permutation-symmetric. No hierarchy.

\textbf{The rotor Hamiltonian and its closed form.} Take $H = \sum_{a<b} \sigma(a,b) \cdot iL_aL_b$, where the couplings $\sigma(a,b)$ are the octonion multiplication-table signs. Independent verification of this work identified, and we confirmed to machine precision, the closed form
\[
H = i(L_u + 2R_u), \qquad u = \sum_{a=1}^{7} e_a,
\]
where $L_u$ and $R_u$ are left and right multiplication by the sum of the imaginary units. Everything that follows is then provable in a few lines (and was confirmed numerically): on $\mathrm{span}\{e_0, u\}$, since $u^2 = -7$, the operator $L_u + 2R_u$ acts as $3M$ with $M^2 = -7\cdot\mathbf{1}$ (a scaled complex structure), giving eigenvalues $\pm 3\sqrt{7}$; on the six-dimensional complement, $R_u = -L_u$ and $L_u^2 = -7\cdot\mathbf{1}$ by alternativity, giving $\pm\sqrt{7}$, each threefold. The unique ground state is
\[
|\Omega_{\mathrm{rotor}}\rangle = \frac{\sqrt{7}\,e_0 - iu}{\sqrt{14}}
\]
--- half identity, half spread uniformly over the seven imaginary units --- with exact invariants: each single-qubit reduced density matrix has spectrum $(1/7, 6/7)$, so $S_1 = S_2 = S_3 = h_2(1/7) \approx 0.5917$; the 3-tangle is $\tau_3 = 8/49$; and all three pairwise concurrences equal $2\sqrt{2}/7$. (The Coffman-Kundu-Wootters identity \cite{coffman} $C_{12}^2 + C_{13}^2 + \tau_3 = 4\lambda_1(1-\lambda_1)$ holds for \emph{every} pure three-qubit state; what is special here is only the equal three-way split, $C_{12}^2 = C_{13}^2 = C_{23}^2 = \tau_3 = 8/49$, which follows from permutation symmetry plus the marginal spectrum.)

\subsection{Controls}
\label{sec:rotorcontrols}

\textbf{Unsigned control.} Replacing the couplings $\sigma(a,b)$ by $+1$ while keeping the true multiplication operators yields a different Hermitian Hamiltonian whose unique ground state is entangled but \emph{hierarchical} and convention-\emph{dependent} (in the Cayley--Dickson convention used here: sorted entropies $(0.577, 0.472, 0.326)$, $\tau_3 = 4/49$ exactly; across the 128 sign gauges the sorted triple takes nine distinct values, unlike the $\sigma$-structured vacuum's single one --- ``unsigned'' is not a basis-independent notion). The table signs are therefore necessary for the symmetric, convention-invariant vacuum --- not for entanglement as such. (An earlier version of this control, like an earlier version of the main computation, suffered from passing a non-Hermitian matrix to a Hermitian eigensolver; both are corrected here, and we flag the failure mode for others automating such computations.)

\textbf{Genericity control.} Random $\mathfrak{so}(7)$ rotor Hamiltonians (generic antisymmetric couplings in place of $\sigma$) also have unique ground states; 16.4\% produce ordered entropies $S_1 > S_2 > S_3$ --- consistent with the $1/6$ exchange-symmetry chance rate --- and 0.9\% of these land in a widened coupling-matching window ($|r_S - 1.8174| < 0.05$; Section~\ref{sec:exp1} uses 0.01), consistent with the Haar-generic rate at that tolerance given the ordering. Generic quadratic dynamics is generic; the $\sigma$-structured choice is distinguished only by its symmetry.

\subsection{Convention Robustness}
\label{sec:robustness}

\textbf{Status: provable from the closed form; confirmed computationally (Experiment 6; \texttt{experiments/results/exp6\_results.json}).}

Because $H = i(L_u + 2R_u)$ with $u$ the sum of the imaginary basis units, the vacuum has the same coordinate vector $(\sqrt{7}, -i, -i, -i, -i, -i, -i, -i)/\sqrt{14}$ in every basis convention: sign gauges $e_a \to s_a e_a$ and relabelings $e_a \to e_{\pi(a)}$ permute and re-sign the summands of $u$ coherently with the basis itself. We nonetheless ran the full battery as a check on the code and the claim: all $2^7 = 128$ sign gauges; all $7! = 5040$ index relabelings (the 168-element XOR-linear $GL(3,2)$ subgroup tracked separately); 2000 random combined transformations; the opposite algebra; right-multiplication operators in place of left; and $H \to -H$. In every case the ground state is unique with sorted entropy triple $(0.5917, 0.5917, 0.5917)$ and $\tau_3 = 8/49$. The invariance is exact, as the closed form requires.

\subsection{What This Does and Does Not Show}
\label{sec:rotorscope}

Because $H = i(L_u + 2R_u)$, the construction is equivalent to distinguishing the unit direction $\hat{u} = u/\sqrt{7}$ in the imaginary octonions, relative to the basis frame that also defines the qubit factorization. The convention-invariance of the vacuum is a provable consequence --- the ``sum of basis units'' direction is permutation- and sign-gauge-covariant by construction --- rather than an empirical discovery. And the choice of $\hat{u}$ matters: other distinguished directions give isospectral Hamiltonians $i(L_{\hat{u}} + 2R_{\hat{u}})$ whose ground states $(e_0 - i\hat{u})/\sqrt{2}$ can be \emph{separable} --- for $\hat{u} = e_1$ this is exactly the preferred-complex-direction state found entropy-degenerate in Section~\ref{sec:octonionic}. So the algebra alone does not select a vacuum; the algebra \emph{plus the frame defining the qubit factorization} selects one. The result is best read as a clean, exactly solvable instance of the same permutation rigidity found kinematically in Section~\ref{sec:octonionic}, now at the dynamical level --- which is exactly the control function this appendix serves: even the canonical dynamical construction yields a symmetric vacuum, not a hierarchy. Within the consistency envelope of Section~\ref{sec:envelope}, a symmetric algebra-derived state is the natural candidate boundary state at a unification scale, and its entropy $0.5917$ lies inside the feasible mid-band --- but as Section~\ref{sec:corollary} records, this compatibility does not constrain the envelope.

\textbf{Relation to prior work.} Multiplication operators on the octonions and the decomposition of $\mathfrak{so}(8)$ into left- and right-multiplication parts are classical material \cite{harvey}, and the $\pm\sqrt{7}$-type spectra of such operators are standard; the three-qubit/octonion/hyperdeterminant nexus is extensively developed in the black-hole/qubit correspondence \cite{borsten}. We have not found this specific Hamiltonian proposed as a three-qubit vacuum selector, or this state's entanglement invariants recorded, but the ingredients are elementary and the result should be regarded as an observation assembled from standard parts, not as new mathematics.

\section*{Note Added}

After this manuscript was completed, a follow-up manuscript made the coupling-entropy question well-posed by replacing single-qubit entropies with entanglement relative to the gauge-invariant operator subalgebras --- the algebraic notion of Barnum-Knill-Ortiz-Viola \cite{barnum} and Zanardi \cite{zanardi}, notion (ii) of Table~\ref{tab:notions} --- and computed the answer exactly in the Jordan--Wigner toy of Section~\ref{sec:jw}, in theorem form. Its Theorem~1 identifies the commutants of the $U(3)$, $SU(3)$, and $U(1)$ gauge actions analytically (by Schur's lemma, with independent numerical certification); its Theorem~2 shows that for every pure state the quantum piece of the algebraic entropy vanishes (``all edge, no bulk''); and its Theorem~3 shows that the weighted-W coupling-matching curve of Section~\ref{sec:wfamily} --- this paper's one surviving static ``match,'' here given in closed form by Lemma~2 (Section~\ref{sec:wfamily}) --- is pure gauge: every state on it has zero gauge-invariant entanglement relative to all three invariant subalgebras, and the W state itself is a gauge transform of a single-particle state. The static match was therefore matching frame artifacts, which strengthens the ill-posedness conclusion of Section~\ref{sec:obstruction} from a diagnosis to a computed fact. The same manuscript also answers the well-posed question in the negative on independent grounds: the gauge-invariant data of the toy reduce to charge-sector statistics with no natural three-fold gauge structure, and --- in its scoped Lemma~4 --- no $SU(2)_L$ action is present within the one-generation, number-conserving three-mode realization, so a three-coupling question was never posable in this minimal setting; a two-ideal/chiral construction is required first.

\begin{thebibliography}{31}

\bibitem{hurwitz} A. Hurwitz, ``\"Uber die Composition der quadratischen Formen,'' Nachr. Ges. Wiss. G\"ottingen, 309 (1898).

\bibitem{gunaydin} M. G\"unaydin and F. G\"ursey, ``Quark structure and the octonions,'' J. Math. Phys. 14, 1651 (1973).

\bibitem{dixon94} G.M. Dixon, \emph{Division Algebras: Octonions, Quaternions, Complex Numbers and the Algebraic Design of Physics}, Springer (1994).

\bibitem{dixon04} G.M. Dixon, ``Division algebras: Family replication,'' J. Math. Phys. 45, 3878 (2004).

\bibitem{furey14} C. Furey, ``Generations: three prints, in colour,'' JHEP 2014, 046 (2014).

\bibitem{furey18} C. Furey, ``$SU(3)_C \times SU(2)_L \times U(1)_Y$ ($\times U(1)_X$) as a symmetry of division algebraic ladder operators,'' Eur. Phys. J. C 78, 375 (2018).

\bibitem{todorov18} I. Todorov and M. Dubois-Violette, ``Deducing the symmetry of the standard model from the automorphism and structure groups of the exceptional Jordan algebra,'' Int. J. Mod. Phys. A 33, 1850118 (2018).

\bibitem{todorov22} I. Todorov, ``Octonion internal space algebra for the Standard Model,'' arXiv:2206.06912 (2022).

\bibitem{szangolies} J. Szangolies, ``The Standard Model Symmetry and Qubit Entanglement,'' Entropy 27(6), 569 (2025); arXiv:2512.17328.

\bibitem{mosseri} R. Mosseri and R. Dandoloff, ``Geometry of entangled states, Bloch spheres and Hopf fibrations,'' J. Phys. A 34, 10243 (2001).

\bibitem{bernevig} B.A. Bernevig and H.-D. Chen, ``Geometry of the three-qubit state, entanglement and division algebras,'' J. Phys. A 36, 8325 (2003).

\bibitem{bekenstein} J.D. Bekenstein, ``Black holes and entropy,'' Phys. Rev. D 7, 2333 (1973).

\bibitem{maldacena} J. Maldacena, ``The large-N limit of superconformal field theories and supergravity,'' Adv. Theor. Math. Phys. 2, 231 (1998).

\bibitem{ryu} S. Ryu and T. Takayanagi, ``Holographic derivation of entanglement entropy from AdS/CFT,'' Phys. Rev. Lett. 96, 181602 (2006).

\bibitem{page} D.N. Page and W.K. Wootters, ``Evolution without evolution,'' Phys. Rev. D 27, 2885 (1983).

\bibitem{takayanagi} T. Takayanagi, ``Emergent Holographic Spacetime from Quantum Information,'' Phys. Rev. Lett. 134, 240001 (2025); arXiv:2506.06595. See also the Simons Foundation ``It from Qubit'' collaboration.

\bibitem{pdg} Particle Data Group, ``Review of Particle Physics,'' Prog. Theor. Exp. Phys. (updated regularly; coupling values used here from the 2022 edition, R.L. Workman et al., 2022, 083C01).

\bibitem{higuchi} A. Higuchi, A. Sudbery, and J. Szulc, ``One-qubit reduced states of a pure many-qubit state: polygon inequalities,'' Phys. Rev. Lett. 90, 107902 (2003).

\bibitem{nielsen} M.A. Nielsen and I.L. Chuang, \emph{Quantum Computation and Quantum Information}, Cambridge University Press (2000).

\bibitem{coffman} V. Coffman, J. Kundu, and W.K. Wootters, ``Distributed entanglement,'' Phys. Rev. A 61, 052306 (2000).

\bibitem{jordanwigner} P. Jordan and E. Wigner, ``\"Uber das Paulische \"Aquivalenzverbot,'' Z. Phys. 47, 631 (1928).

\bibitem{zamolodchikov} A.B. Zamolodchikov, ``Irreversibility of the flux of the renormalization group in a 2D field theory,'' JETP Lett. 43, 730 (1986).

\bibitem{casini} H. Casini and M. Huerta, ``A c-theorem for the entanglement entropy,'' J. Phys. A 40, 7031 (2007).

\bibitem{amaldi} U. Amaldi, W. de Boer, and H. F\"urstenau, ``Comparison of grand unified theories with electroweak and strong coupling constants measured at LEP,'' Phys. Lett. B 260, 447 (1991).

\bibitem{vidal} G. Vidal, ``Entanglement renormalization,'' Phys. Rev. Lett. 99, 220405 (2007).

\bibitem{sawicki} A. Sawicki, M. Walter, and M. Ku\'s, ``When is a pure state of three qubits determined by its single-particle reduced density matrices?,'' J. Phys. A: Math. Theor. 46, 055304 (2013).

\bibitem{borsten} L. Borsten, D. Dahanayake, M.J. Duff, H. Ebrahim, and W. Rubens, ``Black holes, qubits and octonions,'' Phys. Rep. 471, 113 (2009).

\bibitem{stoica} O.C. Stoica, ``Leptons, quarks, and gauge from the complex Clifford algebra $\mathrm{Cl}(6)$,'' Adv. Appl. Clifford Algebras 28, 52 (2018); arXiv:1702.04336.

\bibitem{zanardi} P. Zanardi, ``Virtual quantum subsystems,'' Phys. Rev. Lett. 87, 077901 (2001).

\bibitem{barnum} H. Barnum, E. Knill, G. Ortiz, and L. Viola, ``A subsystem-independent generalization of entanglement,'' Phys. Rev. Lett. 92, 107902 (2004).

\bibitem{harvey} F.R. Harvey, \emph{Spinors and Calibrations}, Academic Press (1990).

\end{thebibliography}

\end{document}
